\section{Introduction}
\label{sec:introduction}

The instruction set architecture of RISC-V has been well accepted in both academic and industrial designs due to its openness and flexibility~\cite{riscv_spec}. PicoRV32 is a small RV32 core specifically targeted for FPGA implementation with low resource usage~\cite{picorv32_yosys}. Although the small core is highly desirable for embedded systems, the arithmetic units in the baseline cores are primarily area-optimized, causing increased latency in arithmetic-intensive execution paths.

This is particularly relevant for edge-AI/TinyML workloads, which operate in close proximity to the sensor. In these cases, compact models still require a large number of integer-based multiply accumulate operations. In these systems, adder critical path and multiplication latency have a strong impact on inference latency as well as energy per inference. Always-on keyword spotting, anomaly detection in industrial sensoring, and low-resolution vision inference are a few examples of the relevant workloads where arithmetic acceleration is known to have a direct impact on system responsiveness and duty-cycled energy efficiency.

In order to mitigate the above limitations, this work proposes the optimization of the PicoRV32 arithmetic datapath. To achieve this, two hardware modifications have been proposed: the first is the integration of a 32-bit Kogge-Stone parallel prefix adder in place of the standard ripple-carry adder, and the second is the integration of a $32\times32$ Vedic multiplier-based PCPI accelerator.

This is because the proposed integration strategy based on the PCPI metric is designed as a modular approach, where the optimized architecture remains as a CPU core-level IP extension, allowing it to be plugged into various SoC scenarios (microcontroller, IoT sensing, or accelerative edge AI subsystems) without any changes to the pipeline.

The paper is prepared based on the existing project outputs as baseline evidence, including architecture, implementation, and results in IEEE conference style.
